\documentclass[aspectratio=169,12pt]{beamer}

\usetheme{Madrid}
\usecolortheme{default}

\usepackage[utf8]{inputenc}
\usepackage[T1]{fontenc}
\usepackage[french]{babel}
\usepackage{listings}
\usepackage{xcolor}
\usepackage{booktabs}
\usepackage{array}
\usepackage{tikz}
\usetikzlibrary{shapes.geometric, arrows.meta, positioning, fit, backgrounds}

% ── Couleurs du thème ─────────────────────────────────────────────────────────
\definecolor{bleuPrimaire}{RGB}{25, 70, 140}
\definecolor{bleuClair}{RGB}{220, 232, 250}
\definecolor{vertCode}{RGB}{40, 120, 60}
\definecolor{codebg}{RGB}{245,245,245}

\setbeamercolor{palette primary}{bg=bleuPrimaire, fg=white}
\setbeamercolor{palette secondary}{bg=bleuPrimaire!80, fg=white}
\setbeamercolor{palette tertiary}{bg=bleuPrimaire!60, fg=white}
\setbeamercolor{frametitle}{bg=bleuPrimaire, fg=white}
\setbeamercolor{block title}{bg=bleuPrimaire, fg=white}
\setbeamercolor{block body}{bg=bleuClair}
\setbeamercolor{alerted text}{fg=red!70!black}

% ── Listings ──────────────────────────────────────────────────────────────────
\lstset{
  backgroundcolor=\color{codebg},
  basicstyle=\ttfamily\tiny,
  keywordstyle=\color{bleuPrimaire}\bfseries,
  commentstyle=\color{vertCode}\itshape,
  stringstyle=\color{red!60!black},
  numbers=none,
  breaklines=true,
  frame=single,
  rulecolor=\color{gray!40},
  showstringspaces=false
}

\lstdefinelanguage{YAML}{
  keywords={true,false,null,on,off,uses,with,run,steps,jobs,name,needs,if,stage,stages,script,image,services,when,branch,only},
  keywordstyle=\color{bleuPrimaire}\bfseries,
  comment=[l]{\#},
  commentstyle=\color{vertCode}\itshape,
  morestring=[b]',
  morestring=[b]",
  stringstyle=\color{red!60!black}
}

\lstdefinelanguage{Groovy}{
  keywords={pipeline,agent,environment,stages,stage,steps,post,when,branch,sh,parallel,always,success,failure,credentials,tools},
  keywordstyle=\color{bleuPrimaire}\bfseries,
  comment=[l]{//},
  commentstyle=\color{vertCode}\itshape,
  morestring=[b]',
  morestring=[b]",
  stringstyle=\color{red!60!black}
}

% ── Informations ──────────────────────────────────────────────────────────────
\title[CI/CD BiblioGest]{Pipeline CI/CD Moderne Multi-Outils}
\subtitle{Application sur le projet BiblioGest}
\author{Génie Logiciel -- Module Gestion des Projets IT}
\date{2025 -- 2026}

% ══════════════════════════════════════════════════════════════════════════════
\begin{document}

% ── Titre ─────────────────────────────────────────────────────────────────────
\begin{frame}
  \titlepage
\end{frame}

% ── Sommaire ──────────────────────────────────────────────────────────────────
\begin{frame}{Sommaire}
  \tableofcontents[hideallsubsections]
\end{frame}

% ══════════════════════════════════════════════════════════════════════════════
\section{Contexte et objectifs}

\begin{frame}{Pourquoi CI/CD~?}
  \begin{columns}
    \begin{column}{0.5\textwidth}
      \textbf{Problèmes sans CI/CD}
      \begin{itemize}
        \item Tests exécutés manuellement, souvent oubliés
        \item Fusions de branches tardives et conflictuelles
        \item Déploiements manuels, sources d'erreurs
        \item Régression détectée en production
      \end{itemize}
    \end{column}
    \begin{column}{0.5\textwidth}
      \textbf{Bénéfices avec CI/CD}
      \begin{itemize}
        \item Feedback immédiat après chaque commit
        \item Qualité mesurée de façon objective
        \item Déploiement reproductible et auditable
        \item Détection des régressions au plus tôt
      \end{itemize}
    \end{column}
  \end{columns}

  \vspace{0.6em}
  \begin{block}{Définitions}
    \textbf{CI} (Intégration Continue)~: automatiser build + tests à chaque push.\\
    \textbf{CD} (Déploiement Continu)~: automatiser jusqu'à la mise en production.
  \end{block}
\end{frame}

\begin{frame}{Objectifs du projet -- Sujet 1}
  \begin{enumerate}
    \item Comparer les principaux moteurs CI/CD~: déclencheurs, runners, jobs,
          stages, artefacts, secrets, notifications
    \item Mettre en place une pipeline complète pour \textbf{BiblioGest}~:
          \begin{itemize}
            \item \textbf{Lint} (Checkstyle + SpotBugs)
            \item Tests unitaires et d'intégration (JUnit~5 + Mockito)
            \item Analyse de qualité (SonarCloud)
            \item Packaging WAR (Maven)
            \item Containerisation Docker + push (Docker Hub + GHCR)
            \item \textbf{Scan sécurité} (Trivy)
            \item \textbf{Notification} Slack
          \end{itemize}
    \item Transposer la même logique sur \textbf{5 moteurs}~: GitHub Actions,
          GitLab CI, Jenkins, CircleCI, Travis CI
    \item Documenter les difficultés réelles rencontrées
  \end{enumerate}

  \vspace{0.4em}
  \alert{Outil principal retenu~: \textbf{GitHub Actions} (minimal, intégré, gratuit OSS)}
\end{frame}

% ══════════════════════════════════════════════════════════════════════════════
\section{Comparaison des outils CI/CD}

\begin{frame}{Concepts communs à tous les outils}
  \begin{columns}
    \begin{column}{0.48\textwidth}
      \begin{description}[\textbf{Notification}]
        \item[\textbf{Trigger}] Événement déclencheur (push, PR, cron)
        \item[\textbf{Runner}] Machine d'exécution (cloud ou self-hosted)
        \item[\textbf{Job}] Unité d'exécution autonome
        \item[\textbf{Stage}] Groupe logique de jobs
      \end{description}
    \end{column}
    \begin{column}{0.48\textwidth}
      \begin{description}[\textbf{Notification}]
        \item[\textbf{Artefact}] Fichier transmis entre jobs (WAR, rapports)
        \item[\textbf{Secret}] Variable sensible chiffrée
        \item[\textbf{Cache}] Répertoire partagé entre runs (Maven, Docker)
        \item[\textbf{Notification}] Message en fin de pipeline (Slack, mail)
      \end{description}
    \end{column}
  \end{columns}
\end{frame}

\begin{frame}{Tableau comparatif des 5 outils}
  \centering
  \scriptsize
  \renewcommand{\arraystretch}{1.3}
  \begin{tabular}{>{\bfseries}p{2.3cm} p{1.9cm} p{1.9cm} p{1.9cm} p{1.9cm} p{1.7cm}}
    \toprule
    Critère & GitHub Actions & GitLab CI/CD & Jenkins & CircleCI & Travis CI \\
    \midrule
    Type             & SaaS              & SaaS + SH     & On-prem        & SaaS        & SaaS \\
    Config           & YAML              & YAML          & Groovy DSL     & YAML        & YAML \\
    Concept central  & jobs / steps      & stages / jobs & pipeline       & workflows   & stages \\
    Marketplace      & 21k+ actions      & Modéré        & 1.9k+ plugins  & Orbs        & Limité \\
    Parallélisme     & \texttt{needs}    & Par stage     & \texttt{parallel} & Natif    & Par stage \\
    Docker natif     & Oui (Buildx)      & Oui (DinD)    & Plugin         & Oui (VM)    & Oui \\
    Gratuit OSS      & \textcolor{vertCode}{Illimité} & Limité & Self-host & 6k min   & \textcolor{red!70!black}{Payant 2021+} \\
    Apprentissage    & \textcolor{vertCode}{Faible} & \textcolor{vertCode}{Faible} & \textcolor{red!70!black}{Élevée} & Moyen & \textcolor{vertCode}{Faible} \\
    \bottomrule
  \end{tabular}
\end{frame}

% ══════════════════════════════════════════════════════════════════════════════
\section{Le projet BiblioGest}

\begin{frame}{BiblioGest -- Architecture}
  \begin{columns}
    \begin{column}{0.45\textwidth}
      \textbf{Stack technique}
      \begin{itemize}
        \item Jakarta EE 10 / Tomcat 10.1
        \item JPA / Hibernate 6 + Oracle XE
        \item JAX-RS (Jersey) -- API REST
        \item JSP + JSTL + Tailwind CSS
        \item Maven pour le build
      \end{itemize}

      \vspace{0.5em}
      \textbf{Route CI/CD testée}\\
      \texttt{GET /api/livres}\\
      \small Retourne le catalogue en JSON,\\
      sert de healthcheck Docker
    \end{column}
    \begin{column}{0.52\textwidth}
      \centering
      \begin{tikzpicture}[
        box/.style={rectangle, draw=bleuPrimaire!50, fill=bleuClair,
                    rounded corners=3pt, minimum width=2.3cm,
                    minimum height=0.75cm, font=\scriptsize\bfseries,
                    text centered},
        arr/.style={-Stealth, thick, bleuPrimaire}
      ]
        \node[box] (nav)  {Navigateur};
        \node[box, below=0.55cm of nav]  (srv) {Servlet / JAX-RS};
        \node[box, below=0.55cm of srv]  (dao) {DAO / JPA};
        \node[box, below=0.55cm of dao]  (db)  {Oracle XE};

        \draw[arr] (nav) -- (srv) node[midway,right]{\tiny HTTP};
        \draw[arr] (srv) -- (dao);
        \draw[arr] (dao) -- (db);
        \draw[arr, dashed] (db)  -- (dao);
        \draw[arr, dashed] (dao) -- (srv);
        \draw[arr, dashed] (srv) -- (nav) node[midway,right]{\tiny JSON};
      \end{tikzpicture}
    \end{column}
  \end{columns}
\end{frame}

% ══════════════════════════════════════════════════════════════════════════════
\section{Architecture de la pipeline}

\begin{frame}{Architecture en 7 étapes}
  \centering
  \begin{tikzpicture}[
    stage/.style={rectangle, draw=bleuPrimaire, fill=bleuClair,
                  rounded corners=4pt, minimum width=1.7cm,
                  minimum height=1cm, font=\scriptsize\bfseries, text centered},
    note/.style={font=\tiny, text=gray},
    arr/.style={-Stealth, thick, bleuPrimaire}
  ]
    \node[stage] (l)  {Lint};
    \node[stage, right=0.45cm of l]  (t) {Tests};
    \node[stage, right=0.45cm of t]  (q) {Qualité};
    \node[stage, right=0.45cm of q]  (p) {Package};
    \node[stage, right=0.45cm of p]  (d) {Docker};
    \node[stage, right=0.45cm of d]  (s) {Sécurité};
    \node[stage, right=0.45cm of s]  (n) {Notify};

    \draw[arr] (l) -- (t);
    \draw[arr] (t) -- (q);
    \draw[arr] (t.south) to[bend right=15] (p.south);
    \draw[arr] (p) -- (d);
    \draw[arr] (d) -- (s);
    \draw[arr] (s) -- (n);

    \node[note, below=0.15cm of q] {main};
    \node[note, below=0.15cm of d] {main};
    \node[note, below=0.15cm of s] {main};
    \node[note, below=0.15cm of n] {main};
  \end{tikzpicture}

  \vspace{0.5em}
  \scriptsize
  \begin{tabular}{>{\bfseries}p{1.6cm} p{8.5cm}}
    \toprule
    Lint     & Checkstyle (style) + SpotBugs (sémantique) \\
    Tests    & Surefire (\texttt{*Test}) + Failsafe (\texttt{*IT}) + JaCoCo \\
    Qualité  & SonarCloud (analyse + quality gate) \\
    Package  & Build du WAR + upload artifact \\
    Docker   & Build multi-stage + push Docker Hub + GHCR \\
    Sécurité & Scan Trivy CVE de l'image \\
    Notify   & Notification Slack succès / échec \\
    \bottomrule
  \end{tabular}
\end{frame}

% ══════════════════════════════════════════════════════════════════════════════
\section{Lint (Checkstyle + SpotBugs)}

\begin{frame}[fragile]{Lint -- Analyse statique du code}
  \begin{columns}
    \begin{column}{0.48\textwidth}
      \textbf{Checkstyle (style)}
      \begin{itemize}
        \item Imports inutilisés
        \item Longueur des lignes
        \item Conventions de nommage
        \item Tabulations interdites
      \end{itemize}

      \vspace{0.5em}
      \textbf{SpotBugs (sémantique)}
      \begin{itemize}
        \item NullPointerException potentielles
        \item Ressources non fermées
        \item Exposition de représentation interne
        \item Sérialisation incorrecte
      \end{itemize}
    \end{column}
    \begin{column}{0.5\textwidth}
      \begin{lstlisting}[language=YAML]
lint:
  runs-on: ubuntu-latest
  steps:
    - uses: actions/checkout@v4
    - uses: actions/setup-java@v4
      with:
        java-version: '17'
        cache: maven
    - run: mvn -B checkstyle:check
    - run: mvn -B compile -DskipTests
    - run: mvn -B spotbugs:check
      continue-on-error: true
    - uses: actions/upload-artifact@v4
      with:
        name: lint-reports
        path: |
          target/checkstyle-result.xml
          target/spotbugsXml.xml
      \end{lstlisting}
    \end{column}
  \end{columns}

  \vspace{0.3em}
  \small \alert{Configuré en mode \textit{informational}~: les violations sont rapportées sans bloquer le pipeline.}
\end{frame}

% ══════════════════════════════════════════════════════════════════════════════
\section{Tests unitaires et d'intégration}

\begin{frame}{Stratégie de test -- contrainte Oracle}
  \begin{alertblock}{Problème}
    Oracle XE met 90s à 3 min à démarrer. Le lancer en CI à chaque push est
    trop lent et trop lourd.
  \end{alertblock}

  \vspace{0.5em}
  \textbf{Solution~: zéro base de données dans les tests}

  \begin{columns}
    \begin{column}{0.5\textwidth}
      \textbf{Tests unitaires (Surefire)}\\[0.3em]
      Classes pures, sans JPA~:
      \begin{itemize}
        \item \texttt{PasswordUtilTest} -- SHA-256
        \item \texttt{BiblioUtilsTest} -- frais, pénalités
      \end{itemize}
    \end{column}
    \begin{column}{0.5\textwidth}
      \textbf{Tests d'intégration (Failsafe)}\\[0.3em]
      DAO mocké avec Mockito~:
      \begin{itemize}
        \item \texttt{LivreResourceIT} -- routes JAX-RS
        \item Codes HTTP 200, 404
        \item Filtrage \texttt{disponible}
      \end{itemize}
    \end{column}
  \end{columns}

  \vspace{0.5em}
  \small Convention Maven~: \texttt{*Test.java} → Surefire \quad|\quad
  \texttt{*IT.java} → Failsafe
\end{frame}

\begin{frame}[fragile]{Exemple -- Test unitaire et d'intégration}
  \begin{lstlisting}[language=Java]
// BiblioUtilsTest.java
@Test
void penaliteTroisJoursLesMillerables() {
    BigDecimal p = BiblioUtils.calculerPenalite(
            LocalDate.of(2026, 4, 19),
            LocalDate.of(2026, 4, 22),
            new BigDecimal("5.00"));
    assertEquals(new BigDecimal("22.50"), p);  // 3j * 5.00 * 1.5
}
  \end{lstlisting}

  \begin{lstlisting}[language=Java]
// LivreResourceIT.java - test d'integration
@ExtendWith(MockitoExtension.class)
class LivreResourceIT {
    @Mock LivreDao dao;
    LivreResource resource;

    @BeforeEach void setUp() { resource = new LivreResource(dao); }

    @Test
    void getByIdInexistantRetourne404() {
        when(dao.findByIdWithRelations(99L)).thenReturn(null);
        assertEquals(404, resource.getById(99L).getStatus());
    }
}
  \end{lstlisting}
\end{frame}

% ══════════════════════════════════════════════════════════════════════════════
\section{Qualité avec SonarCloud et SonarQube}

\begin{frame}{SonarCloud vs SonarQube self-hosted}
  \centering
  \footnotesize
  \renewcommand{\arraystretch}{1.4}
  \begin{tabular}{>{\bfseries}p{2.8cm} p{4.5cm} p{4.5cm}}
    \toprule
    Critère & SonarCloud (SaaS) & SonarQube self-hosted \\
    \midrule
    Hébergement      & Cloud Sonar     & À votre charge \\
    Coût             & Gratuit OSS public & Community gratuit \\
    Maintenance      & Aucune          & Mises à jour, sauvegardes \\
    Confidentialité  & Code envoyé chez Sonar & 100\% on-prem \\
    Quality gates    & Préconfigurés   & Customisables \\
    Idéal pour       & Démos, OSS      & Entreprises régulées \\
    Setup            & Quelques minutes & Plusieurs heures (infra) \\
    \bottomrule
  \end{tabular}

  \vspace{0.6em}
  \begin{block}{Démonstration SonarQube self-hosted}
    \texttt{docker compose -f docker-compose.sonarqube.yml up -d}\\[0.2em]
    Stack SonarQube 10.4 Community + PostgreSQL 15, accessible sur
    \texttt{http://localhost:9000}
  \end{block}
\end{frame}

\begin{frame}[fragile]{Intégration SonarCloud dans le pipeline}
  \begin{lstlisting}[language=YAML]
quality:
  name: SonarCloud
  runs-on: ubuntu-latest
  needs: test
  if: github.ref == 'refs/heads/main'
  steps:
    - uses: actions/checkout@v4
      with: { fetch-depth: 0 }
    - uses: actions/setup-java@v4
      with: { java-version: '17', distribution: 'temurin', cache: maven }
    - name: Cache SonarCloud
      uses: actions/cache@v4
      with:
        path: ~/.sonar/cache
        key: ${{ runner.os }}-sonar
    - name: Analyse SonarCloud
      env:
        SONAR_TOKEN: ${{ secrets.SONAR_TOKEN }}
      run: |
        mvn -B verify sonar:sonar \
          -Dsonar.projectKey=${{ vars.SONAR_PROJECT_KEY }} \
          -Dsonar.organization=${{ vars.SONAR_ORGANIZATION }} \
          -Dsonar.host.url=https://sonarcloud.io
  \end{lstlisting}

  \small Couverture JaCoCo lue automatiquement via
  \texttt{sonar.coverage.jacoco.xmlReportPaths}.
\end{frame}

% ══════════════════════════════════════════════════════════════════════════════
\section{Containerisation Docker}

\begin{frame}[fragile]{Dockerfile multi-stage}
  \begin{columns}
    \begin{column}{0.55\textwidth}
      \begin{lstlisting}
# Stage 1: Build Maven
FROM eclipse-temurin:17-jdk-alpine AS builder
WORKDIR /build
COPY pom.xml ./
RUN apk add --no-cache maven && \
    mvn dependency:go-offline || true
COPY src ./src
COPY checkstyle.xml ./
RUN mvn package -DskipTests

# Stage 2: Runtime Tomcat
FROM tomcat:10.1-jdk17-temurin-alpine AS runtime
RUN rm -rf /usr/local/tomcat/webapps/*
COPY --from=builder /build/target/bibliotheque.war \
     /usr/local/tomcat/webapps/ROOT.war
ENV DB_URL=jdbc:oracle:thin:@oracle-xe:1521/XEPDB1
RUN addgroup -S tomcat && adduser -S tomcat -G tomcat && \
    chown -R tomcat:tomcat /usr/local/tomcat
USER tomcat
EXPOSE 8080
HEALTHCHECK --start-period=90s \
  CMD wget -qO- http://localhost:8080/api/livres || exit 1
CMD ["catalina.sh", "run"]
      \end{lstlisting}
    \end{column}
    \begin{column}{0.42\textwidth}
      \textbf{Avantages}
      \begin{itemize}
        \item L'image finale ne contient pas Maven
        \item Pas de code source dans l'image
        \item Surface d'attaque réduite
        \item Image \textasciitilde 250 Mo au lieu de \textasciitilde 700 Mo
      \end{itemize}

      \vspace{0.5em}
      \textbf{Sécurité runtime}
      \begin{itemize}
        \item Exécution sous user non-root \texttt{tomcat}
        \item Healthcheck sur \texttt{/api/livres}
        \item Labels OCI standard
      \end{itemize}
    \end{column}
  \end{columns}
\end{frame}

\begin{frame}[fragile]{Job Docker -- multi-registry (Docker Hub + GHCR)}
  \begin{lstlisting}[language=YAML]
docker:
  needs: package
  if: github.ref == 'refs/heads/main'
  permissions:
    contents: read
    packages: write
  steps:
    - uses: actions/checkout@v4
    - uses: docker/setup-buildx-action@v3
    - uses: docker/login-action@v3
      with: { username: ${{ secrets.DOCKER_USERNAME }}, password: ${{ secrets.DOCKER_PASSWORD }} }
    - uses: docker/login-action@v3
      with: { registry: ghcr.io, username: ${{ github.actor }}, password: ${{ secrets.GITHUB_TOKEN }} }
    - id: meta
      uses: docker/metadata-action@v5
      with:
        images: |
          ${{ secrets.DOCKER_USERNAME }}/bibliotheque
          ghcr.io/${{ github.repository_owner }}/bibliotheque
        tags: |
          type=raw,value=latest
          type=sha,format=long
    - uses: docker/build-push-action@v5
      with: { context: ., push: true, tags: ${{ steps.meta.outputs.tags }},
              cache-from: type=gha, cache-to: type=gha,mode=max }
  \end{lstlisting}
\end{frame}

% ══════════════════════════════════════════════════════════════════════════════
\section{Sécurité Trivy et notification Slack}

\begin{frame}[fragile]{Scan sécurité Trivy}
  \begin{columns}
    \begin{column}{0.5\textwidth}
      \textbf{Que détecte Trivy~?}
      \begin{itemize}
        \item CVE des paquets OS (Alpine, Debian)
        \item CVE des dépendances Maven/npm/etc.
        \item Mauvaises configurations
        \item Secrets exposés
      \end{itemize}

      \vspace{0.5em}
      \textbf{Format SARIF}
      \begin{itemize}
        \item Standard de l'industrie
        \item Affiché dans l'onglet \emph{Security} de GitHub
        \item Permet de suivre les CVE dans le temps
      \end{itemize}
    \end{column}
    \begin{column}{0.48\textwidth}
      \begin{lstlisting}[language=YAML]
security-scan:
  needs: docker
  if: github.ref == 'refs/heads/main'
  steps:
    - uses: aquasecurity/trivy-action@0.24.0
      with:
        image-ref: ${{ secrets.DOCKER_USERNAME }}/bibliotheque:latest
        severity: 'CRITICAL,HIGH'
        ignore-unfixed: true
    - uses: aquasecurity/trivy-action@0.24.0
      with:
        image-ref: ${{ secrets.DOCKER_USERNAME }}/bibliotheque:latest
        format: 'sarif'
        output: 'trivy-results.sarif'
    - uses: github/codeql-action/upload-sarif@v3
      with:
        sarif_file: 'trivy-results.sarif'
      \end{lstlisting}
    \end{column}
  \end{columns}
\end{frame}

\begin{frame}[fragile]{Notification Slack en fin de pipeline}
  \begin{lstlisting}[language=YAML]
notify:
  needs: [lint, test, quality, package, docker, security-scan]
  if: always() && github.ref == 'refs/heads/main'
  steps:
    - id: status
      run: |
        if [ "${{ contains(needs.*.result, 'failure') }}" = "true" ]; then
          echo "text=Pipeline ECHOUEE" >> $GITHUB_OUTPUT
          echo "color=danger" >> $GITHUB_OUTPUT
        else
          echo "text=Pipeline REUSSIE" >> $GITHUB_OUTPUT
          echo "color=good" >> $GITHUB_OUTPUT
        fi
    - uses: slackapi/slack-github-action@v1.27.0
      with:
        payload: |
          {
            "attachments": [{
              "color": "${{ steps.status.outputs.color }}",
              "blocks": [ ... champs branche/commit/auteur/lien ... ]
            }]
          }
      env:
        SLACK_WEBHOOK_URL: ${{ secrets.SLACK_WEBHOOK_URL }}
  \end{lstlisting}

  \small \alert{\texttt{if: always()}} garantit l'envoi même en cas d'échec
  d'un job précédent.
\end{frame}

% ══════════════════════════════════════════════════════════════════════════════
\section{Équivalences multi-outils}

\begin{frame}[fragile]{GitLab CI/CD -- Équivalent}
  \begin{lstlisting}[language=YAML]
stages: [lint, test, quality, package, docker, security, notify]

variables:
  MAVEN_OPTS: "-Dmaven.repo.local=$CI_PROJECT_DIR/.m2/repository"

.maven-cache: &maven-cache
  cache:
    key: "$CI_COMMIT_REF_SLUG-maven"
    paths: [.m2/repository]

test:
  stage: test
  <<: *maven-cache
  image: eclipse-temurin:17-jdk
  script: [mvn $MAVEN_CLI_OPTS verify]
  artifacts:
    reports:
      junit: [target/surefire-reports/TEST-*.xml, target/failsafe-reports/TEST-*.xml]
  coverage: '/Total.*?([0-9]{1,3})%/'

docker:build-push:
  stage: docker
  image: docker:24
  services: [docker:24-dind]
  rules: [{ if: '$CI_COMMIT_BRANCH == "main"' }]
  script:
    - docker login -u $DOCKER_USERNAME -p $DOCKER_PASSWORD
    - docker build -t $DOCKER_USERNAME/bibliotheque:$CI_COMMIT_SHA .
    - docker push  $DOCKER_USERNAME/bibliotheque:$CI_COMMIT_SHA
  \end{lstlisting}

  \small Différence clé~: jobs d'un même stage en \textbf{parallèle automatique},
  sans \texttt{needs}.
\end{frame}

\begin{frame}[fragile]{Jenkins -- Jenkinsfile déclaratif}
  \begin{lstlisting}[language=Groovy]
pipeline {
    agent any
    environment {
        SONAR_TOKEN  = credentials('sonar-token')
        DOCKER_CREDS = credentials('docker-hub')
    }
    tools { jdk 'jdk-17'; maven 'maven-3.9' }
    stages {
        stage('Lint') {
            parallel {
                stage('Checkstyle') { steps { sh 'mvn -B checkstyle:check' } }
                stage('SpotBugs')   { steps { sh 'mvn -B compile -DskipTests'
                                              sh 'mvn -B spotbugs:check || true' } }
            }
        }
        stage('Tests') { steps { sh 'mvn -B verify' }
            post { always { junit 'target/**/TEST-*.xml' } } }
        stage('Docker') { when { branch 'main' }
            steps {
                sh 'docker build -t bibliotheque:${GIT_COMMIT} .'
                sh 'docker push ${DOCKER_CREDS_USR}/bibliotheque:latest'
            } }
    }
    post {
        success { slackNotify('good',   'REUSSIE') }
        failure { slackNotify('danger', 'ECHOUEE') }
    }
}
  \end{lstlisting}

  \small Force de Jenkins~: \texttt{parallel} et \texttt{post} natifs, mais
  Groovy plus complexe que YAML.
\end{frame}

\begin{frame}[fragile]{CircleCI -- Orbs et workflows}
  \begin{lstlisting}[language=YAML]
version: 2.1

orbs:
  maven:  circleci/maven@1.4.1
  docker: circleci/docker@2.6.0

executors:
  java-executor:
    docker: [{ image: cimg/openjdk:17.0 }]
    resource_class: medium

jobs:
  test:
    executor: java-executor
    steps:
      - checkout
      - restore_cache: { keys: ["maven-{{ checksum \"pom.xml\" }}"] }
      - run: mvn -B verify --no-transfer-progress
      - save_cache: { key: "maven-{{ checksum \"pom.xml\" }}", paths: [~/.m2] }
      - store_test_results: { path: target/surefire-reports }

workflows:
  bibliogest:
    jobs:
      - lint
      - test: { requires: [lint] }
      - docker-build-push: { requires: [test], filters: { branches: { only: main } } }
  \end{lstlisting}

  \small Spécificités~: \textbf{orbs} (modules réutilisables), \textbf{executors} nommés,
  \textbf{workspaces} pour partager des artefacts.
\end{frame}

\begin{frame}[fragile]{Travis CI -- Syntaxe minimaliste}
  \begin{lstlisting}[language=YAML]
language: java
jdk: openjdk17
services: [docker]
cache: { directories: [$HOME/.m2] }

stages:
  - lint
  - test
  - name: docker
    if: branch = main

jobs:
  include:
    - stage: lint
      script:
        - mvn -B checkstyle:check
        - mvn -B compile -DskipTests
        - mvn -B spotbugs:check || true
    - stage: test
      script: mvn -B verify --no-transfer-progress
    - stage: docker
      script:
        - echo "$DOCKER_PASSWORD" | docker login -u "$DOCKER_USERNAME" --password-stdin
        - docker build -t $DOCKER_USERNAME/bibliotheque:$TRAVIS_COMMIT .
        - docker push  $DOCKER_USERNAME/bibliotheque:$TRAVIS_COMMIT

notifications:
  email: { on_success: change, on_failure: always }
  \end{lstlisting}

  \small Avantage~: la syntaxe la plus \textbf{concise} des cinq.
  Inconvénient~: \alert{plan gratuit OSS supprimé en 2021, communauté en fuite}.
\end{frame}

% ══════════════════════════════════════════════════════════════════════════════
\section{Difficultés rencontrées}

\begin{frame}{Difficultés rencontrées}
  \begin{enumerate}
    \item \textbf{Tests sans base de données}\\
          Le bloc \texttt{static} de \texttt{JpaUtil} échoue si Oracle est absent.
          Solution~: classes utilitaires pures + injection de mocks via constructeur.

    \vspace{0.3em}
    \item \textbf{Constructeur d'injection package-private}\\
          \texttt{LivreResource(LivreDao)} était inaccessible depuis le package
          \texttt{integration}. Correction~: passage en \texttt{public}.

    \vspace{0.3em}
    \item \textbf{Checkstyle bloquant sur code legacy (177 violations)}\\
          Choix pédagogique~: mode \emph{informational}
          (\texttt{failOnViolation=false}). Les rapports sont uploadés comme
          artefacts, sans bloquer le pipeline.

    \vspace{0.3em}
    \item \textbf{SpotBugs et entités JPA}\\
          Les warnings \texttt{EI\_EXPOSE\_REP} sont structurellement nécessaires
          pour Hibernate. \texttt{continue-on-error: true} dans GitHub Actions.

    \vspace{0.3em}
    \item \textbf{Permissions GHCR}\\
          Push sur \texttt{ghcr.io} exige \texttt{permissions: packages: write}
          explicite, sinon erreur 403 obscure.

    \vspace{0.3em}
    \item \textbf{Conflits JSTL / EL dans WEB-INF/lib}\\
          Ajout explicite de \texttt{jakarta.servlet.jsp.jstl-api:3.0.0} en
          \texttt{compile} + \texttt{jakarta.el} passé en \texttt{provided}.

    \vspace{0.3em}
    \item \textbf{Démarrage lent d'Oracle XE en Docker}\\
          90s à 3 min. Résolu par \texttt{condition: service\_healthy} dans
          docker-compose.
  \end{enumerate}
\end{frame}

% ══════════════════════════════════════════════════════════════════════════════
\section{Comparaison finale}

\begin{frame}{Recommandations par contexte}
  \centering
  \footnotesize
  \renewcommand{\arraystretch}{1.5}
  \begin{tabular}{>{\bfseries}p{4.5cm} p{8cm}}
    \toprule
    Contexte & Recommandation \\
    \midrule
    Projet académique / OSS sur GitHub & \textbf{GitHub Actions} -- gratuit, immédiat, marketplace énorme \\
    Équipe sur GitLab self-hosted     & \textbf{GitLab CI/CD} -- cohérence outillage DevOps complet \\
    Grande entreprise on-premise      & \textbf{Jenkins} -- maturité, plugins métier, AD/LDAP \\
    Projet cloud-native polyglotte    & \textbf{CircleCI} -- orbs, performance, multi-VCS \\
    Projet OSS hérité                 & Migrer depuis \textbf{Travis CI} vers GitHub Actions \\
    \bottomrule
  \end{tabular}

  \vspace{0.8em}
  \begin{block}{Pour BiblioGest}
    \textbf{GitHub Actions} a été retenu~: configuration YAML lisible,
    intégration Docker Hub + GHCR + SonarCloud + Slack sans plugin
    supplémentaire, runners gratuits illimités pour OSS public.
  \end{block}
\end{frame}

% ── Conclusion ────────────────────────────────────────────────────────────────
\begin{frame}{Conclusion -- 3 leçons retenues}
  \begin{block}{1. La testabilité conditionne la CI/CD}
    Classes utilitaires pures, injection de dépendances, pas d'initialisation
    DB dans les blocs statiques. L'architecture détermine la rapidité du
    pipeline.
  \end{block}

  \begin{block}{2. Portabilité conceptuelle des pipelines}
    Lint~$\to$~test~$\to$~qualité~$\to$~package~$\to$~docker~$\to$~sécurité~$\to$~notify
    s'exprime sur les 5 moteurs avec des syntaxes différentes mais des
    concepts identiques.
  \end{block}

  \begin{alertblock}{3. Le pipeline n'est pas une fin en soi}
    Un pipeline bien configuré dont les alertes Slack sont ignorées, dont les
    rapports Sonar ne sont jamais consultés, n'offre qu'une \textbf{illusion
    de rigueur}.
  \end{alertblock}

  \vspace{0.5em}
  \centering
  \large\textbf{Merci -- Questions~?}
\end{frame}

\end{document}
